\section{Confiança no resultado}
\label{sec:validacao-quantitativa}

O comportamento descrito foi verificado por dois caminhos independentes.

\paragraph{Verificação analítica.} A distribuição inicial calculada pelo modelo
(Figura~\ref{fig:dist-inicial}) reproduz a solução clássica \(\sinh(\adist(1-x))/\sinh(\adist)\) da
Equação~\eqref{eq:sinh} com erro desprezível para malhas refinadas, confirmando que a capacitância
série foi incorporada corretamente.

\paragraph{Verificação por simulação independente (ATP/EMTP).} O mesmo caso — escada de \Nsec{} seções,
neutro aterrado, com \Cg{} e \Cs{} (\(\adist=5\)) — foi montado e resolvido no ATP, simulador de
transitórios eletromagnéticos consolidado \cite{dommel1969}. A Figura~\ref{fig:atp-comparacao} sobrepõe,
nó a nó, a resposta do solver em Python e a do ATP: as curvas são praticamente indistinguíveis. Os picos
por nó concordam dentro de \(+0{,}003\,\%\) (diferença sistemática que vem apenas da normalização da
fonte) e o erro pontual máximo é \(\le 0{,}011\,\%\) de \SI{1}{kV} na janela de interesse dielétrico
(\SIrange{0}{30}{\micro\second}).

\staticfigure{0.84}{comparacao_atp_gerador.png}{Validação cruzada Python\,\(\times\)\,ATP para o caso do
gerador (neutro aterrado, com capacitância série, \(\adist=5\)), nos nós a 0\,\%, 25\,\%, 50\,\%,
75\,\% e 100\,\% do enrolamento. As curvas dos dois solvers se sobrepõem; o nó a 100\,\% permanece em
\SI{0}{V} (neutro aterrado).}{fig:atp-comparacao}

As duas verificações sustentam que as animações e os perfis deste relatório representam a física real do
enrolamento, e não um artefato do método numérico.
